\section{3. Pureza dos materiais}

\begin{frame}{A pureza que o silício precisa ter}
  \begin{columns}[T]
    \begin{column}{0.52\linewidth}
      \begin{itemize}
        \item O silício da eletrônica tem \alert{99,999999999\%} de pureza,
              os \termo{onze noves}.
        \item Isso é menos de \termo{um átomo estranho por bilhão}.
        \item É um grão de açúcar dissolvido em uma piscina olímpica.
      \end{itemize}
    \end{column}
    \begin{column}{0.44\linewidth}
      \small
      \begin{table}
        \centering
        \begin{tabular}{ll}
          \toprule
          \textbf{Uso} & \textbf{Pureza} \\
          \midrule
          Liga metálica & 98\% \\
          Painel solar  & 99,9999\% \\
          Eletrônica    & 11 noves \\
          \bottomrule
        \end{tabular}
      \end{table}
    \end{column}
  \end{columns}
  \vspace{0.4cm}
  \centering\small
  A dopagem usa um átomo em um milhão. Se a sujeira estiver nessa mesma
  proporção, ela toma o lugar do dopante.
\end{frame}

\begin{frame}{Como o silício é purificado}
  \begin{enumerate}
    \item \termo{Forno}: areia mais carbono a \grau{2000}. Sai um silício com
          98\% de pureza.
    \item \termo{Destilação}: o silício vira gás, o gás é destilado e depois
          volta a ser sólido, já limpo.
    \item \termo{Cristalização}: o cristal é puxado devagar do silício
          derretido, e a sujeira fica no líquido.
  \end{enumerate}
  \vspace{0.4cm}
  \gancho{No amplificador de uma antena, impureza vira ruído, e mais ruído
  significa menos alcance para a estação.}
\end{frame}
